\section*{第 \ref{ch:b3:total-synthesis} 章——全合成：策略与经典路线}

\begin{solution}{exo:b3:total-synthesis:1}
在 90\,\% 时：$0.9^{10} = 35$\,\%，$0.9^{20} = 12$\,\%。在 80\,\% 时：$0.8^{10} = 11$\,\%，$0.8^{20} = 1.2$\,\%。
\end{solution}

\begin{solution}{exo:b3:total-synthesis:2}
共 $5 + 3 + 1 + 4 = 13$ 步；\omterm{def:b3:total-synthesis:architecture}{最长线性步骤} $5 + 1 + 4 = 10$。
\end{solution}

\begin{solution}{exo:b3:total-synthesis:3}
4-（二甲氨基）丁-2-酮，$\ce{CH3COCH2CH2N(CH3)2}$：丙酮的烯醇加到亚胺离子
$\ce{CH2=N+(CH3)2}$ 上。
\end{solution}

\begin{solution}{exo:b3:total-synthesis:4}
切断与 C=C 双键相邻的两个环键，得到丁-1,3-二烯和丁-3-烯-2-酮
（甲基乙烯基酮），二者加热即可反应。
\end{solution}

\begin{solution}{exo:b3:total-synthesis:5}
线性路线：$0.9^{21} = 11$\,\%。汇聚路线：$0.9^{11} = 31$\,\%，几乎是前者的三倍。
\end{solution}

\begin{solution}{exo:b3:total-synthesis:6}
A：$(7 + 1)/12 = 67$\,\%。B：$7/9 = 78$\,\%。
\end{solution}

\begin{solution}{exo:b3:total-synthesis:7}
不加这四步时：$0.9^8 = 43$\,\%。若加上四个收率均为 95\,\% 的额外步骤：$43\,\% \times 0.95^4 = 35$\,\%。
\end{solution}

\begin{solution}{exo:b3:total-synthesis:8}
伯醇：用体积较大的硅醚（\textit{叔}-丁基二苯基硅基）保护，用氟离子脱除。
仲醇：用 4-甲氧基苄醚保护，用 DDQ 氧化脱除。叔醇
因位阻较大，通常可以不加保护，也可保护为体积较小的硅醚，
用温和酸脱除。每种条件都不影响其他基团。
\end{solution}

\begin{solution}{exo:b3:total-synthesis:9}
早期的\omterm{def:b3:total-synthesis:total}{全合成}需要数十步，\omterm{def:b3:total-synthesis:architecture}{总产率}远低于 1\,\%；
如今可从欧洲紫杉的针叶中提取出一种化合物，它已具备完整的环系和大部分手性碳原子，
欧洲紫杉是可再生资源，只需几步即可转化为目标分子。
\end{solution}

\begin{solution}{exo:b3:total-synthesis:10}
迈克尔加成：二酮的烯醇负离子（C2 位，位于两个羰基之间）加到甲基乙烯基酮的
$\beta$-碳上，生成三酮。羟醛缩合：甲基酮的烯醇负离子
分子内进攻环上的羰基，再脱水得到
环己\omterm{def:b3:radicals-carbenes:wolff}{烯酮}，即维兰德–米舍尔酮，一种带有角
甲基的双环烯二酮。
\end{solution}

\begin{solution}{exo:b3:total-synthesis:11}
$y^{m+1}/y^{2m+1} = y^{-m}$：在 90\,\% 时，$m = 5$ 对应 1.7，$m = 15$ 对应 4.9。路线越长，
汇聚策略的优势越明显。
\end{solution}

\begin{solution}{exo:b3:total-synthesis:12}
每个环都在类椅式构象中由碳正离子与进攻的双键闭合；
每个双键上的加成都是反式的，取代基占据假平伏位置，
从而使每个环稠合均为\textit{反式}。\textit{E} 型双键
使其两侧链的延伸方向恰好得到\textit{反式}稠合；\textit{Z} 型则
得到\textit{顺式}稠合：多烯的几何构型决定了
产物的立体化学。
\end{solution}

\begin{solution}{pb:b3:total-synthesis:1}
\textbf{1.} $\ce{C4H6O2}$，$\ce{CH5N}$，$\ce{C5H6O5}$，$\ce{C8H13NO}$。

\textbf{2.} $\ce{C4H6O2 + CH5N + C5H6O5 -> C8H13NO + 2CO2 + 2H2O}$。

\textbf{3.} 86.0、31.0 和 \qty{146.0}{g/mol}；托品酮为 \qty{139.0}{g/mol}。

\textbf{4.} $139.0/(86.0 + 31.0 + 146.0) = 0.53$：53\,\%。

\textbf{5.} 其中心的 $\ce{CH2}$ 基团位于两个羰基之间，在
水中易于烯醇化；丙酮本身亲核性太弱，活化用的羧基
在反应后离去。

\textbf{6.} 亚胺离子需要游离的胺（酸过多会使胺质子化），而
烯醇化与亚胺离子的生成都需要一定量的酸：pH 接近中性时二者可兼顾。

\textbf{7.} $\ce{CH3NH2}$ 加到其中一个醛基上；失水后得到亚胺离子
$\ce{R-CH=N+H-CH3}$。

\textbf{8.} 丙酮二羧酸的一个烯醇碳原子进攻亚胺离子的碳原子，形成一个
C--C 键并生成仲胺。

\textbf{9.} 仲胺与同一链上的第二个醛基缩合，
生成环状（五元环）亚胺离子。

\textbf{10.} 酮另一侧的烯醇进攻亚胺离子，闭合
六元环，形成双环骨架。

\textbf{11.} 两者均为 $\beta$-酮酸，通过六元环状过渡态失去 $\ce{CO2}$，
生成酮的烯醇。

\textbf{12.} 两个 C--C 键和两个 C--N 键。

\textbf{13.} $0.75^{14} = 0.018$：1.8\,\%。

\textbf{14.} $42/1.78 = 24$。

\textbf{15.} 100\,\%：一步完成，全部为成键步骤。

\textbf{16.} 一步操作，以水为溶剂，原料廉价，无需保护基，
除二氧化碳和水外几乎没有废物。

\textbf{17.} 植物由氨基酸衍生的
环状亚胺离子与乙酸衍生的单元，经同类曼尼希化学构建
托烷骨架。

\textbf{18.} \omterm{def:b3:total-synthesis:cascade}{串联反应}（以及\omterm{def:b3:pericyclic:families}{环加成}等成环反应，如第尔斯–阿尔德
反应）。

\textbf{19.} 否：\omterm{def:b3:point-groups:mirror}{镜面}穿过 N 原子、其甲基、C3 和羰基
氧原子。

\textbf{20.} 每个都有四种不同的基团，但二者在
分子自身的对称面两侧互为镜像：属于内消旋型化合物。

\textbf{21.} 非对映异构体，区别在于 C3 上的 OH 相对于
氮桥的取向不同。

\textbf{22.} 否：两者都保留了\omterm{def:b3:point-groups:mirror}{镜面}。

\textbf{23.} 羰基的两个面不等价：一个面朝向
氮桥，另一个面朝向双碳桥，二者受阻程度不同。

\textbf{24.} 一锅法路线的\omterm{def:b3:total-synthesis:architecture}{总产率}约为 14 步线性路线的
24 倍。
\end{solution}
